% titel: Blatt 8, A1: Siedepunkte ordnen und Kräfte nennen
% Teil II: der durchdachte Weg.

\subsection*{Gegeben und gesucht}
BaCl$_2$, HF, CO und Ne sollen nach ihrem Siedepunkt geordnet werden; die
wirksamen intermolekularen Kräfte sind zu nennen.

\subsection*{Ansatz}
Art der Teilchen und der zugehörigen Anziehungskräfte bestimmen. Ionen
$>$ H-Brücke $>$ Dipol-Dipol $>$ Dispersion bei ähnlicher Grösse.

\subsection*{Rechnung und Begründung}
\begin{itemize}
  \item \textbf{BaCl$_2$} ist ein Salzkristall aus Ba$^{2+}$ und 2\,Cl$^-$:
        starke Coulomb-Kräfte $\Rightarrow$ höchster Siedepunkt.
  \item \textbf{HF}: starke Wasserstoffbrücken (F sehr elektronegativ)
        $\Rightarrow$ deutlich über Dipol-Dipol-Stoffen, trotz kleiner Molmasse.
  \item \textbf{CO} ist ein polares kleines Molekül: nur Dipol-Dipol- und
        Dispersion, keine H-Brücke.
  \item \textbf{Ne} ist ein Edelgasatom: nur schwache Dispersion
        $\Rightarrow$ tiefster Siedepunkt.
\end{itemize}

\begin{ergebnis}
Siedepunkte (aufsteigend):
\[
  \text{Ne} < \text{CO} < \text{HF} < \text{BaCl}_2 .
\]
Kräfte: Ne nur Dispersion; CO vor allem Dipol-Dipol (+ Dispersion); HF
Wasserstoffbrücken; BaCl$_2$ ionisch.
\end{ergebnis}

\begin{kontrolle}
Plausibel: Edelgase sieden am tiefsten, ionische Feststoffe am höchsten.
HF liegt klar über CO trotz ähnlicher Molmasse --- genau das macht die
Sonderstellung der H-Brücke sichtbar.
\end{kontrolle}

\begin{merksatz}
Siedepunkt reiht sich mit der Stärke der intermolekularen Kräfte: ionisch
$>$ Wasserstoffbrücke $>$ Dipol-Dipol $>$ Dispersion. Die Molmasse allein
reicht nicht, um die Reihenfolge zu erklären.
\end{merksatz}
